\section{Benchmark Tasks \& Toolkit}\label{sec:tasks}

Every robot learning pipeline makes implicit perception choices that are rarely justified.
RoboVLMs~\citep{li2025robovlms} confirmed that backbone selection dominates VLA success, yet evaluated only in simulation.
\coolname{} provides the real-world evaluation by framing the benchmark as \textbf{five perception decisions} (Table~\ref{tab:task-decisions}), all evaluated \textit{zero-shot} on the same physical scenes and objects.

\begin{table}[t]
\centering
\caption{Five perception decisions. Each evaluates a signal consumed by deployed robot systems.}
\label{tab:task-decisions}
\resizebox{\linewidth}{!}{%
\begin{tabular}{p{0.4cm} p{2.0cm} p{3.5cm} p{2.5cm} p{2.5cm} p{2.5cm}}
\toprule
\textbf{\#} & \textbf{Decision} & \textbf{Systems} & \textbf{Models} & \textbf{Metrics} & \textbf{Diagnostic} \\
\midrule
D1 & Depth model
   & DP3~\citep{ze2024dp3}, GR00T~\citep{bjorck2025gr00t}
   & DA-V2, Video DA, Depth Pro, UniDepth, Metric3D, ZoeDepth
   & AbsRel, RMSE, $\delta\!<\!1.25$
   & STR, TS, SGC \\
\addlinespace
D2 & Detector / grounder
   & OK-Robot~\citep{liu2024okrobot}
   & \textit{Text:} GroundingDINO, OWL-ViT, YOLO-World. \textit{Visual:} DE-ViT, NIDS-Net
   & AP$_{50}$, Acc@0.5
   & STR, text vs.\ visual \\
\addlinespace
D3 & Tracker
   & Pick-and-place, re-grasping
   & SAM2 video, Cutie
   & MOTA, IDF1
   & ID switches/phase \\
\addlinespace
D4 & Scene QA VLM
   & $\pi_0$~\citep{black2024pi0}, OpenVLA~\citep{kim2024openvla}
   & GPT-4o, Qwen2-VL, LLaVA, Molmo~2, PaliGemma
   & Fuzzy-match Acc
   & STR on QA \\
\addlinespace
D5 & Spatial QA VLM
   & Spatial manipulation
   & Same as D4
   & Fuzzy-match Acc
   & STR on spatial QA \\
\bottomrule
\end{tabular}%
}
\end{table}

\paragraph{D1: Monocular depth.}
Predict a dense metric depth map from a single RGB frame (no scale alignment).
Metric depth is the primary input for point-cloud-based grasping~\citep{ze2024dp3} and the pseudo-GT that world models~\citep{zhen2025tesseract} train on.
GT: D435 depth with validity mask.

\paragraph{D2: Detection, grounding, and in-context visual recognition.}
\label{sec:d2-detection-grounding}
Localise target objects from a \textit{text} prompt (category name or referring expression from questionnaire) or a \textit{reference image} (from single-object scene).
\coolname{} supports both modes on the same physical objects, enabling a controlled comparison: \textit{which prompting mode degrades less across phases?}
GT: bounding boxes from instance masks.

\paragraph{D3: Multi-object tracking.}
Maintain consistent object identities across frames.
The interaction phase is the tracking stress test: hand occlusion causes ID switches---a primary failure mode in demonstration data.
GT: SAM2 video-propagated tracklets with human-corrected keyframes; initialised with GT first-frame masks (isolating tracking from detection).
To avoid circularity when evaluating SAM2 as a tracker, we report SAM2's results alongside a note that its GT was generated by the same model family, and we verify that human corrections are sufficient to make the GT independent of any single model's failure modes (Appendix~\ref{sec:appendix-annotation}).

\paragraph{D4: Scene understanding QA.}
Given an RGB frame and a question (``What objects are on the table?''), produce a textual answer.
VLAs are built on VLMs; if the VLM cannot answer correctly on D435 imagery, the VLA is likely to misidentify targets.
Evaluation uses fuzzy matching (normalised string similarity $\geq 0.8$ after lowercasing and article removal) to handle synonym and formatting variation.
Examples: ``a red bottle'' matches ``red bottle'' (\cmark); ``there are 3 items'' matches ``3 items'' (\cmark); ``the bottle is on the table'' does not match ``bottle'' (\xmark, similarity $< 0.8$).
GT: human-authored question templates grounded in robotics operator needs (object identity, attribute verification, count, state change), instantiated deterministically per frame using questionnaire attributes and per-frame visibility masks.
The same question has \textit{different correct answers} across phases.
Unlike static VQA benchmarks (VQAv2, GQA) or 3D-situated QA (SQA3D~\citep{ma2022sqa3d}), \coolname{} tests on robot-grade sensor imagery ($640\!\times\!480$, D435 colour) with phase-varying GT and crowd-sourced natural-language object descriptions.
Template design and answer-key validation details are in Appendix~\ref{sec:appendix-qa}.

\paragraph{D5: Spatial reasoning QA.}
Given a spatial question (``What is to the left of the bottle?''), produce a textual answer.
All spatial relations are defined in the \textit{image plane}, matching how VLMs process images.
Critically, spatial relations \textit{change} across phases---making this, to our knowledge, the first perception benchmark with phase-varying spatial QA ground truth.
GT: human-designed spatial-relation templates (left/right/above/below/near/far), instantiated deterministically from bounding-box geometry; templates were authored to reflect the spatial queries a robot operator or planner would issue.

\paragraph{Toolkit.}
\label{sec:toolkit}
\texttt{rpx-benchmark} (pip-installable, MIT) encapsulates the full pipeline:
\texttt{rpx bench monocular\_depth --hf-checkpoint my-model --split hard}.
Three plugin registries (task, metric, model adapter) enable community extension without core changes.
138 CI-verified tests ensure reproducibility.
Extended documentation in Appendix~\ref{sec:appendix-toolkit}.
